\section{Discussion}\label{sec:discussion}

\subsection{Why temperatures, prices, regularizers, and attention weights are the same kind of object}

The unifying claim of this paper is not that many disciplines happen to reuse the rhetoric of cost. It is that a common mathematical role keeps reappearing whenever a descriptive layer must remain closed under bounded spending. Temperature, market price, regularization strength, and attention-like priority variables all function as dual coordinates of constrained maintenance. They are not interchangeable in substance, but they are the same kind of object in form: each is a shadow price attached to a resource that the layer cannot spend without limit \cite{jaynes1957information,arrowdebreu1954,sims2003rationalinattention}.

That is why the paper insists on a strict, layer-relative definition of currency. What matters is not the physical or semantic meaning of the quantity in isolation, but whether it structurally governs feasibility, protocol accounting, monotone potential, or dual exchange. Under that criterion, energy and budget are primal currencies, while temperature and price are dual currencies; in learning, complexity penalties and their multipliers play the same role; in bounded-representation settings, priority or attention parameters can be interpreted in the same formal register when they price limited bandwidth or precision. The recurrence is therefore not an analogy imported after the fact. It is the repeated appearance of the same closure geometry under different substrates.

\subsection{Why this is a cross-layer story rather than a remark about multipliers}

The familiar statement that Lagrange multipliers exist is correct but incomplete \cite{boyd2004convex}. It describes the local geometry of a constrained optimization problem once the variables, state space, and constraints have already been chosen. The present paper addresses an earlier and more structural question: why do such constrained problems appear across layers in the first place?

The answer proposed here is that packaging and maintenance force them to appear. A higher-level object can persist only if lower-level spending remains bounded while that object is maintained. That bounded spending becomes a feasibility law at the higher layer, and the marginal value of relaxing it becomes the higher-layer currency. The empirical budget sweep makes this point visible. Price is not inserted interpretively on top of the problem. It emerges when the inherited lower-layer budget binds, and it vanishes when the budget becomes slack.

This is also why the paper's main principle is directional. Lower-layer currency becomes higher-layer constraint; higher-layer constraint becomes higher-layer shadow price. The layers are linked by maintenance, not merely by formal duality. That is the sense in which the paper offers an economics of closure rather than a relabeling of existing optimization theory.

\subsection{Why proxy currencies fail, and why cycle-space can appear before strong affinity norms}

The proxy-ablation result sharpens the distinction between structural currencies and convenient correlates. A proxy can predict some behavior without being the quantity the layer is truly spending. But once the task is to recover a stable dual variable and to generalize out of sample, structural alignment matters. The good cost in the ablation inherits its meaning from the lower-layer ledger; the bad cost preserves scale while destroying that alignment. The result is exactly what the strict definition predicts: worse held-out performance and a less stable inferred price. Calling every useful statistic a currency would blur precisely the distinction the experiment makes visible.

The resolution-ladder result adds a complementary lesson. Finer resolution reveals a larger cycle space before it necessarily reveals large surviving affinity norms. That is not a defect of the theory; it is an important structural fact. Resolution can expose the capacity for independent currency coordinates before those coordinates carry substantial active drive at the observed scale. In other words, the dimension of the available currency space can grow before the realized spending in that space becomes large. This is one reason the paper separates the existence of a currency direction from the magnitude currently expressed along it.

\subsection{What the experiments establish, and what they do not}

The empirical results establish five things within a controlled finite-state laboratory. First, directionality audits are honest under coarse observation: pushed-forward trajectories do not manufacture path-reversal asymmetry. Second, shadow price emerges monotonically from active budget and collapses in the slack regime. Third, the number of independent currency directions grows with resolution as loop structure becomes visible. Fourth, structurally aligned currencies outperform proxy costs both predictively and inferentially. Fifth, stronger budget enforcement improves packaging coherence as measured by idempotence defect.

These results do not establish that every real scientific system admits a one-dimensional price description, nor that every cross-layer transition is captured by the simple MaxCal family used here. The paper offers a finite-state realization and a reusable measurement discipline, not a universal reduction. The currency notion is explicitly layer-relative: what counts as a currency depends on the layer's objects, protocols, and audits. The proxy-ablation study is controlled and synthetic, designed to isolate structural alignment rather than to emulate a naturally occurring benchmark. The Lean theorem proves deterministic-pushforward monotonicity for a finite KL quantity in the finite-\(\mathrm{klFun}\) form; it is a formal anchor for the audit side of the argument, not a full mechanization of the currency--constraint principle or of the empirical Markov laboratory.

Those limitations matter because they mark the paper's scope honestly. The contribution is not to have completed a general theory of all cross-scale economics. It is to have identified a precise structural pattern, defined it tightly enough to distinguish genuine currencies from proxies, and demonstrated it in a reproducible finite setting with a partial formal guarantee on the monotonicity side.

\subsection{What this paper adds to the Six Birds series}

Within the Six Birds series, the foundations paper established the primitive vocabulary of closure and made audits, packaging, and definability explicit at the layer level \cite{tsiokos2026sixbirds}. PICA then supplied an interaction algebra and a stochastic laboratory in which closure primitives can be enabled, coupled, and measured \cite{tsiokos2026create}. The present paper adds the spending law of that program. It identifies what a closure must budget to remain stable, shows how that budget becomes the next layer's feasibility boundary, and explains why the next layer responds with a price.

In that sense, this paper sits between theory and measurement. It extends the conceptual side of Six Birds by defining currencies strictly and cross-layerly, and it extends the empirical side by turning that definition into a measurable set of signatures. The result is not just another example paper in the series. It is the point at which closure acquires an economy.
